% Chapter 6 -- Rubric: PO11, PO9 (9.1--9.4), PO6 (6.1), PO7 (7.1), PO8 (8.1--8.4)
\chapter{Project Management, Teamwork, and Impact}
\label{ch:management}

\section{Project Management and Finance}

\subsection{Planning and Risk Management}
\label{sec:plan}

The project ran across one academic session with five contributors, all carrying
a full course load. Work was partitioned along the subsystem boundaries of
\Cref{fig:architecture}, which was chosen partly for architectural reasons and
partly because those boundaries are the ones that let several people work
without blocking each other: the engine, the content pipeline and the client
each depend on the ontology's \emph{interface} but not on each other's internals.

\Cref{fig:gantt} shows the schedule. The version-control history records 93
commits between April and September 2026, and the phases below are drawn from
it rather than from an idealised plan.

\begin{figure}[htbp]
  \centering
  \begin{ganttchart}[
      hgrid, vgrid, x unit=0.78cm,
      time slot format=isodate-yearmonth, time slot unit=month,
      title label font=\small, bar label font=\small,
      bar/.style={fill=gray!40}
    ]{2026-03}{2026-12}
    \gantttitlecalendar{year, month=shortname} \\
    \ganttbar{Requirements \& research}{2026-03}{2026-04} \\
    \ganttbar{Engine design \& BKT core}{2026-04}{2026-06} \\
    \ganttbar{Ontology v1}{2026-05}{2026-07} \\
    \ganttbar{API \& frontend}{2026-06}{2026-08} \\
    \ganttbar{Content pipeline}{2026-07}{2026-09} \\
    \ganttbar{Ontology rebuild}{2026-09}{2026-09} \\
    \ganttbar{Testing \& hardening}{2026-08}{2026-10} \\
    \ganttbar{Report \& presentation}{2026-10}{2026-12}
  \end{ganttchart}
  \caption{Project schedule. The ontology rebuild in September was not in the
    original plan; it was forced by Investigation 2 (\Cref{sec:investigation}).}
  \label{fig:gantt}
\end{figure}

\paragraph{Where actual progress differed from plan.} Three deviations, each
instructive.

\begin{enumerate}
  \item \textbf{A late, unplanned ontology rebuild.} The plan assumed the
    knowledge structure was a one-time input. It turned out to be the project's
    critical path: the first ontology was too sparse for the gating policy to
    act on, and rebuilding it consumed most of September. The schedule absorbed
    this only because the compiler-based build made a rebuild a one-command
    operation.
  \item \textbf{A discarded content-generation approach.} Weeks of work on local
    model inference were written off when measurement showed it could not finish
    in the remaining schedule (Investigation 1). Migrating to a hosted API
    recovered the capability but consumed schedule that had been allocated to
    coverage.
  \item \textbf{Scope reduction on authentication.} Deferred deliberately once
    it became clear the ontology work would not otherwise finish.
\end{enumerate}

\begin{table}[htbp]
  \centering
  \caption{Risk register: principal risks, their realisation, and mitigation}
  \label{tab:risks}
  \begin{tabularx}{\textwidth}{>{\hsize=0.75\hsize}Y>{\hsize=0.4\hsize}Y>{\hsize=1.85\hsize}Y}
    \toprule
    Risk & Realised? & Mitigation and outcome \\
    \midrule
    Content generation infeasible at required scale & \textbf{Yes} & Measured early rather than assumed; switched to a hosted API with pooled keys. Cost the schedule, not the deliverable. \\
    Knowledge structure inadequate for the inference built on it & \textbf{Yes} & Detected by structural analysis rather than by a user complaint; rebuild was affordable because the ontology build was automated from the start. \\
    Machine-generated questions are wrong or off-topic & Partly & Two-model generate-then-verify with blind re-solve; only items passing both checks are written. \\
    Source corpus corruption propagating silently & \textbf{Yes} & One malformed record destroyed 334 neighbours before the parser was made resynchronising; escape repair made unconditional. \\
    External API rate limits or cost overrun & \textbf{Yes} & Key pooling, adaptive backoff, batch verification, resumable progress files. One parallelised run was halted mid-flight when its consumption rate proved unsustainable, and the work was restructured to a cheaper sequential method. \\
    Knowledge loss across a five-person, multi-month project & Partly & A chain of handover documents recording state and rationale per subsystem; these were written precisely because context kept being lost between sessions and contributors. \\
    Team member unavailability during examinations & Partly & Subsystem boundaries chosen so that no single absence blocked others. \\
    Secret leakage into a public repository & No & Credentials confined to a git-ignored environment file from the first commit that needed one; never committed. \\
    \bottomrule
  \end{tabularx}
\end{table}

\subsection{Budgeting and Resource Identification}
\label{sec:budget}

The project was built entirely on free tiers and hardware the team already
owned; direct monetary expenditure was zero. That is not the same as costless,
so \Cref{tab:budget} itemises both the actual outlay and the notional value of
what was consumed, and \Cref{tab:opex} models what a real deployment would cost.

\begin{table}[htbp]
  \centering
  \caption{Development resources: actual and notional cost}
  \label{tab:budget}
  \begin{tabularx}{\textwidth}{YY>{\raggedleft\arraybackslash}p{0.18\textwidth}}
    \toprule
    Resource & Basis & Actual cost \\
    \midrule
    Engineering effort & 5 students, one academic session, part-time alongside coursework & BDT 0 (in kind) \\
    Development hardware & Personal laptops, CPU only, no GPU purchased & BDT 0 (owned) \\
    Database and hosting & Supabase free tier & BDT 0 \\
    Model API for generation & Gemini free-tier keys, pooled across accounts to raise the rate ceiling & BDT 0 \\
    Embedding model & Open-weight multilingual model, run locally & BDT 0 \\
    Source content & Textbooks and question banks already owned by team members & BDT 0 \\
    Version control and CI & GitHub free tier & BDT 0 \\
    \midrule
    \textbf{Total direct expenditure} & & \textbf{BDT 0} \\
    \bottomrule
  \end{tabularx}
\end{table}

The free tiers are a genuine constraint rather than a lucky break: the pipeline's
key-pooling, adaptive backoff and resumable-progress design exist specifically
because per-key rate limits, not money, were the scarce resource.

\begin{table}[htbp]
  \centering
  \caption{Indicative operating cost for a production deployment (order-of-magnitude estimate, not a quotation)}
  \label{tab:opex}
  \begin{tabularx}{\textwidth}{>{\hsize=0.8\hsize}Y>{\hsize=1.5\hsize}Y>{\hsize=0.7\hsize}Y}
    \toprule
    Item & Driver & Scale \\
    \midrule
    One-time question generation & $1{,}688$ skills $\times$ 6 Bloom levels $= 10{,}128$ specifications, each a generation call plus a share of a batched verification call & Dominant one-time cost \\
    Managed database & Tens of MB of content; one row per learner per skill & Lowest paid tier adequate at pilot scale \\
    Application hosting & One small always-on process; inference is constant-time per response & Smallest instance class \\
    Content refresh & Regeneration only when the syllabus or ontology changes & Infrequent \\
    Domain expert review & Chemist and physicist review of the ontology (\Cref{sec:eval-limits}) & The largest genuinely unavoidable human cost \\
    \bottomrule
  \end{tabularx}
\end{table}

% TODO (optional): if you can get real BDT figures for the paid tiers you would
% use, put them in Table 6.3 -- an examiner reads a costed table more
% favourably than a qualitative one.

\subsection{Economic Analysis and Financial Prospecting}
\label{sec:economics}

The cost structure has one useful property: it is almost entirely fixed. Content
generation and ontology construction are one-time costs; serving an additional
learner costs a database row and a few constant-time updates. The marginal cost
per learner is therefore close to zero, which is the necessary condition for any
viable offering in this market, because \Cref{sec:market} established that
individual willingness to pay for software is low even though willingness to pay
for preparation is high.

Three revenue paths follow from that, in descending order of realism:

\begin{enumerate}
  \item \textbf{Institutional licensing.} Coaching centres and schools pay per
    cohort for diagnostic reporting that they cannot produce manually. The buyer
    has budget, the value proposition is diagnosis at a granularity a human
    cannot reach across a large class, and the number of buyers is small enough
    to reach directly.
  \item \textbf{Freemium to learners.} Diagnostic free, unlimited practice paid.
    This fits observed behaviour but yields little per user and depends on scale
    the project does not have.
  \item \textbf{Public or NGO deployment.} Offered at no cost as an equity
    intervention (\Cref{sec:equity}), funded by a body whose objective is access
    rather than revenue. This is the path most consistent with the system's
    social value, and the least dependent on a commercial market.
\end{enumerate}

The honest financial reading is that this is not, in its current form, a
business: it is a prototype whose economics work only after the one-time content
cost is paid and whose evidence of efficacy --- the thing an institution would
actually buy --- does not yet exist (\Cref{sec:eval-limits}). Establishing that
evidence is a prerequisite for any of the three paths above, not a follow-on.

\subsection{Sustainability in Business and Commercialisation}
\label{sec:commercialisation}

For the system to outlive the course, three things must be true. First, the
knowledge structure must be maintainable by someone who is not the original
team: this is addressed architecturally, since the ontology is compiled from a
source file by a documented one-command build, with editorial decisions
isolated in a single configuration file. Second, the content must be
regenerable rather than hand-curated, which the pipeline provides. Third, the
system must be transferable --- which is why the repository carries an
operational guide and a chain of handover documents written for a reader with no
prior context, and why the test harness runs without credentials.

The realistic continuation path is academic rather than commercial: the
codebase, the validated ontology and the documented open questions form a
defined starting point for a subsequent capstone or thesis, with the domain
expert review and a controlled efficacy study as the natural next pieces of
work. A commercial path would require a sponsor to fund the generation run and
the expert review before any user-facing launch.

\section{Individual and Teamwork}
\label{sec:team}

\subsection{Participation and Contribution}
\label{sec:participation}

% TODO: This is the one section that cannot be written from the repository.
% The table below records the subsystem ownership that IS visible in the git
% history (93 commits, five contributors) -- map each row to the right member,
% correct anything misattributed, and add each member's own account of
% decision-making participation and on-time completion. An examiner will expect
% names here, not roles alone.

\begin{table}[htbp]
  \centering
  \caption{Contribution by subsystem}
  \label{tab:contributions}
  \begin{tabularx}{\textwidth}{>{\hsize=0.5\hsize}Y>{\hsize=0.6\hsize}Y>{\hsize=1.9\hsize}Y}
    \toprule
    Member & Primary area & Contributions \\
    \midrule
    Member A (ID) & Adaptive engine & BKT implementation, the three-phase mastery policy, the skill-graph data structure with its cycle check, and the policy specification document that the code is audited against. \\
    Member B (ID) & Backend API \& sessions & FastAPI service, session lifecycle for diagnostic and practice, the tiered question search, and the spillover rules. \\
    Member C (ID) & Content pipeline & Corpus parsing and repair, retrieval index, the generate-then-verify generation pipeline, and the database loader with its validation pass. \\
    Member D (ID) & Frontend & React client, the question-answer flow, mastery visualisation, bilingual interface, and mathematical rendering. \\
    Member E (ID) & Ontology & Skill extraction from the corpus, the ontology compiler, the full-corpus rebuild and its editorial pass, and the structural validators. \\
    \bottomrule
  \end{tabularx}
\end{table}

Ownership was primary rather than exclusive --- the ontology rebuild in
particular drew in contributors from outside its owning area, because it sat on
every other subsystem's critical path.

\subsection{Collaboration and Conflict Resolution}
\label{sec:collaboration}

Collaboration was organised around the repository. Work proceeded on
subsystem-aligned branches merged into a main line, so that the integration
point was explicit and reviewable rather than continuous. The credential-free
test harness served as the shared definition of ``not broken'': any contributor
could run it in seconds before merging, which removed most integration
disagreements by making them checkable rather than arguable.

The substantive technical disagreements in the project were about the ontology
--- specifically about granularity, since finer skills improve diagnosis but
multiply content cost, and different contributors weighted those differently.
These were resolved by writing the criteria down and testing them against data:
the decision to rebuild was settled by the structural analysis in Investigation
2 rather than by argument, and the granularity target was fixed by reference to
what the gating policy needs to function. The general pattern the team
converged on was to convert a disagreement about judgement into a measurement
wherever one could be defined.

% TODO: add one or two concrete instances -- a specific disagreement, how it was
% raised, and how it was settled. Named specifics read as genuine; general
% statements read as filler.

\subsection{Leadership and Direction}
\label{sec:leadership}

Direction was distributed by subsystem rather than centralised: each area's
owner set its technical direction and was accountable for it, with cross-cutting
decisions --- the model family, the granularity target, the decision to rebuild
--- taken jointly. The mechanism that kept direction coherent across a
five-and-a-half-month project was documentary: each significant change was
accompanied by a handover note recording not just what changed but why, so that
a contributor returning to an area after weeks away, or a new contributor
entering it, inherited the reasoning rather than only the code.

The clearest instance of planning after a setback is Investigation 1. When local
generation was measured and found unable to finish within the schedule, the
response was not to persist but to re-plan: the approach was abandoned, the
alternative adopted, and the reasoning recorded so that the abandoned path would
not be re-attempted by someone who had not seen the measurements.

% TODO: add how the team kept members motivated through the September crunch and
% the examination period -- an examiner is looking for evidence of the human
% mechanics here, not only the technical ones.

\subsection{Multidisciplinary Engagement}
\label{sec:multidisciplinary}

The project is multidisciplinary by construction: it required computer science
(inference, graph algorithms, systems), educational psychology (knowledge
tracing, Bloom's taxonomy, the zone of proximal development), and subject-matter
expertise in three natural sciences at HSC level. The team's engagement outside
its own discipline took three forms: reading and applying the education research
literature rather than treating pedagogy as common sense; working directly with
HSC Mathematics, Physics and Chemistry source material to build the ontology;
and explicitly identifying the boundary of its own competence --- the decision to
mark 32 graph fragments and a body of authored chemistry edges as requiring a
chemist's and a physicist's review, rather than resolving them by a
computer scientist's intuition, is the clearest example of recognising where
another discipline's authority was needed.

% TODO: if any teacher, domain expert, or student outside the team reviewed
% content or tested the system, name that engagement here -- it strengthens this
% indicator considerably and should also appear in the Acknowledgement.

\section{Impact on Society}
\label{sec:society}

\textbf{Positive: access.} The dominant positive impact is distributional. Expert
diagnostic attention is currently rationed by money and geography --- the best
coaching is concentrated in cities and priced accordingly. A system whose
marginal cost per learner is near zero can offer a form of that attention to a
student in a district town with a phone and a connection. This is the project's
principal social justification.

\textbf{Positive: efficiency of effort.} Redirecting a student from repeated
failure on a downstream topic to the upstream skill actually missing reduces
wasted study time in a year where time is the binding constraint.

\textbf{Negative: misdiagnosis.} The corresponding harm is specific and worth
stating plainly. A wrong prerequisite edge or a miscalibrated parameter can send
a learner to study something they already know, or release them from something
they do not. In a high-stakes year that is a real cost borne by someone with no
way to detect the error. The design response is the conservative mastery
threshold and the refusal to guess prerequisite edges mechanically
(Investigation 3); the honest response is that this risk is why efficacy
evidence is a precondition for deployment, not a nicety.

\textbf{Negative: intensification.} A system that always has a next question can
extend an already unhealthy preparation culture. The design does not currently
counter this --- there are no session limits or break prompts --- and adding them
is listed in future work.

\textbf{Health and safety.} No physical risk. The relevant health dimension is
psychological: adaptive systems that target a constant challenge level
deliberately keep the learner near their failure boundary, which is
motivationally corrosive if presented as bare failure. Mitigation in the current
design is that mastery is shown as a graph of what has been secured, not only as
what remains.

\textbf{Cultural.} Bangla-first content with correct mathematical typesetting is
a deliberate cultural position: it refuses the common assumption that serious
technical study must happen in English, and it serves students whose
mathematical vocabulary is Bangla.

\textbf{Legal.} Two obligations govern deployment: personal data of a
predominantly under-20 user base, and copyright in source material. Both are
addressed in \Cref{sec:standards}; the current prototype's lack of
authentication means it is not yet lawful to deploy publicly with real learner
data, which is stated as a blocking condition rather than a limitation.

\section{Environmental Impact and Life Cycle Sustainability}
\label{sec:sustainability}

\textbf{Development.} Modest: five laptops over one session, and CPU-bound
embedding of a corpus of a few thousand records. The abandoned local-inference
approach would have been materially worse --- weeks of continuous CPU
computation --- so the decision taken on schedule grounds in Investigation 1 also
reduced energy consumption.

\textbf{Content generation.} This is the project's environmental hotspot. Large
language model inference carries a real energy and water footprint in the
provider's data centre, and the specification space at rebuild scale is 10{,}128
generation calls plus batched verification. Three design choices reduce it, all
originally made for cost or throughput reasons: verification uses a smaller
model than generation; verification calls are batched across several
specifications; and progress is checkpointed so an interrupted run resumes
rather than restarting. A fourth mitigation is intrinsic --- generation is a
one-time cost, not a per-learner one, and it amortises over every learner
thereafter.

\textbf{Operation.} Very low. One small always-on process and a managed
database; inference per response is arithmetic, not model inference. The system
does no work while a learner is reading a question.

\textbf{Retirement.} Two obligations. Learner data must be deletable, which the
schema supports through cascading deletion from the learner record, and which a
production deployment should expose as an explicit account-deletion path. The
generated question bank and ontology are the project's durable artefacts and are
worth preserving beyond the system itself --- as a syllabus-aligned skill graph
they have value independent of this application, and releasing them rather than
discarding them is the sustainable end-of-life choice.

\section{Ethics}
\label{sec:ethics}

\subsection{Equity and Inclusivity}
\label{sec:equity}

\textbf{Language.} The system is Bangla-first for content, with a bilingual
interface. This is the single largest inclusivity decision in the project: it
admits students whose technical vocabulary is Bangla, who are systematically
disadvantaged by English-only tools.

\textbf{Cost.} Zero marginal cost per learner makes free provision technically
possible, which is the precondition for reaching students for whom coaching fees
are prohibitive.

\textbf{Access assumptions, and who is still excluded.} The design assumes a
browser and a connection. That excludes students without either --- which is
precisely the group the access argument above claims to serve. This tension is
not resolved in the current system. Offline capability and low-bandwidth
operation are the honest mitigations, and neither is implemented.

\textbf{Disability.} Mathematics is rendered as accessible markup rather than as
images, and controls are keyboard-operable, but no formal accessibility audit
against WCAG AA \cite{wcag21} was performed and no screen-reader testing was
done. Stating this as an outstanding gap is more useful than claiming compliance
that has not been verified.

\textbf{Fairness across learners.} One model-level equity issue deserves naming:
the parameters in \Cref{tab:bloom-params} are uniform across all learners. A
learner whose true guess or slip behaviour differs from those priors --- for
instance, a careful learner who slips rarely --- is modelled slightly wrongly, and
uniformly. Individualised parameters \cite{yudelson2013individualized} are the
known remedy and require response data the system does not yet have.

\subsection{Accountability and Personal Responsibility}
\label{sec:accountability}

Responsibility for the system's behaviour rests with the team, and within the
team it is traceable: subsystem ownership is recorded
(\Cref{tab:contributions}) and every change is attributable through version
control. The project's accountability practices were concrete rather than
declarative --- the defect in \Cref{sec:revisions} was reported, analysed and
guarded against rather than quietly fixed, and the two over-flagging validators
in Investigation 4 were documented as team errors of judgement rather than
deleted.

The accountability commitment that matters most for this system is about claims.
The team is responsible for not asserting that ATLAS improves examination
outcomes, because no evidence for that exists. That commitment is why
\Cref{sec:eval-limits} exists in this report in the form it does.

\subsection{Proper Use of Intellectual Property}
\label{sec:ip}

\textbf{Prior art and literature.} The model family, the taxonomy and the
retrieval approach are not original to this project and are cited to their
sources \cite{corbett1995knowledge,bloom1956taxonomy,anderson2001revised,doignon1985spaces,lewis2020retrieval,vygotsky1978mind}.
The original contribution claimed is the composition --- Bloom-conditioned BKT
with ancestor pull-up and conjunctive gating over a syllabus-derived graph --- not
its components.

\textbf{Third-party software.} All third-party components are open-source under
permissive licences and are used as dependencies rather than copied into the
codebase: FastAPI, Uvicorn, Pydantic (backend); React, Vite, Tailwind CSS,
KaTeX, Lucide (frontend); sentence-transformers and NumPy (pipeline); PostgreSQL
via Supabase. Their licences are carried in the dependency manifests.
% TODO: if the department requires an explicit licence table, expand this into
% one listing each package and its licence (MIT / Apache-2.0 / BSD).

\textbf{Source content.} The corpus is drawn from copyrighted HSC textbooks and
published question banks. It is used as a private retrieval and grounding
corpus; it is not redistributed, is excluded from distributable artefacts, and
generated items are new text rather than reproductions. Any public deployment
would require this position to be reviewed with the rights holders.

\textbf{Generative AI disclosure.} Generative AI was used substantively in this
project and is disclosed accordingly:
\begin{itemize}
  \item \textbf{As a system component.} The question bank is machine-generated
    via a hosted large language model, with a second model verifying each item.
    This is not an incidental use --- it is a designed part of the architecture
    (\Cref{sec:generation-pipeline}), and the verification stage exists
    precisely because unverified model output is not trustworthy as assessment
    content.
  \item \textbf{As a development tool.} AI coding assistants were used during
    development for code generation, refactoring, documentation, the ontology
    extraction and consolidation passes, and in the preparation of this report.
    All such output was reviewed by team members before being committed, and the
    team accepts full responsibility for it (\Cref{sec:accountability}).
    Investigation 4 is a direct record of that review process catching
    machine-produced errors --- both over-flagging validators were AI-assisted
    artefacts corrected by hand review.
\end{itemize}
% TODO: adjust the second bullet to match exactly what your department requires
% to be disclosed, and how specific it expects you to be about which tools.

\subsection{Professionalism and Ethical Codes}
\label{sec:codes}

The ACM Code of Ethics \cite{acm2018code} and the IEEE Code of Ethics
\cite{ieee2020code} both place avoiding harm and honest representation of work
above the practitioner's own interest. Three project decisions follow directly
from those obligations and were in tension with the team's interest in
presenting a polished result:

\begin{enumerate}
  \item \textbf{Refusing an unevidenced efficacy claim} (ACM 1.3, ``be honest and
    trustworthy''). The straightforward way to strengthen this report would be
    to assert that the system improves learning. \Cref{sec:eval-limits} instead
    states that no such evidence was gathered.
  \item \textbf{Declining to connect prerequisite edges mechanically}
    (ACM 1.2, ``avoid harm''). Accepting all 83 machine proposals would have
    produced a better-looking connectivity statistic. Thirty-five were rejected
    because a wrong edge silently corrupts a learner's diagnosis, which is a harm
    the statistic would have concealed.
  \item \textbf{Documenting defects rather than quietly closing them}
    (IEEE 7.8.7, accepting and offering honest criticism of technical work). The
    topic-code defect, the corpus corruption and the over-flagging validators are
    all recorded in the repository and in this report with their causes.
\end{enumerate}

The public-interest dimension specific to this system is that its users are
young people making consequential decisions under pressure, largely without the
expertise to audit the advice they are given. That asymmetry is the reason the
team treats ``do not claim more than the evidence supports'' as the project's
governing professional obligation rather than as a reporting formality.
